\documentclass[10pt,a4paper]{amsart}
\usepackage[margin=19mm]{geometry}
\usepackage{amsmath,amssymb,mathtools}
\usepackage[scr=rsfs,cal=euler]{mathalfa}
\usepackage{xfrac}
\usepackage[unicode,hidelinks,bookmarks=false]{hyperref}
\newcommand{\ie}{i.e.\ }
\newcommand{\cf}{cf.\ }
\newcommand{\resp}{respectively\ }
\newcommand{\Subsection}{\subsection}
\providecommand{\nobreakdash}{\nobreak\hbox{-}}
\setlength{\emergencystretch}{2em}
\multlinegap0pt
\usepackage{bm}
\usepackage{bbm}
\usepackage{dsfont}
\usepackage{xspace}
\usepackage{cancel}
\usepackage{mathtools}
\usepackage{amsthm}
\makeatletter
\@ifundefined{Th}{\newtheorem{Th}{Theorem}}{}
\makeatother
\newcommand{\wT}{\widetilde{T}}
\title[Around the Merino--Welsh conjecture: improving Jackson's ine]{Around the Merino--Welsh conjecture: improving Jackson's inequality}
\author{P\'eter Csikv\'ari}
\date{Source: arXiv:2502.19196v2 (CC BY 4.0)}
\begin{document}
\maketitle

\noindent\emph{Source and licence.} Around the Merino--Welsh conjecture: improving Jackson's inequality. Version arXiv:2502.19196v2 (CC BY 4.0), distributed under the Creative Commons Attribution 4.0 International licence, as recorded in the arXiv licence metadata for this exact version and shown on the abstract page https://arxiv.org/abs/2502.19196v2. Full rights notice: licenses/arxiv-2502.19196-CC-BY-4.0.txt.\par\smallskip

\noindent\textit{Editorial scope.} Counted unit: the Th whose source label is \texttt{None} in the article (source0.tex lines 618--627, source label \texttt{None}), delivered together with the article's own complete original proof of it (source0.tex lines 629--691, one \texttt{proof} environment of 63 lines). No mathematics is written, repaired or re-derived: statement and proof are the article's own text line for line. Required context retained uncounted: eq-vlsz (lines 446--448). Every definition, display and label of those lines is retained unchanged. Omitted from this package and not redistributed: 1--445, 449--617, 628--628, 692--1326. Nothing inside a retained excerpt is omitted or altered: every retained body is delivered exactly as the article's lines are written, so no omission receipt of any kind accompanies this package. Citations: the retained excerpts carry no citation command at all, so no bibliography entry of the article is transcribed and no reference is redistributed. Reference adaptation: none was needed and none was made; every reference inside the delivered unit resolves inside it, so no unresolved reference or citation is produced. Printed slips kept literal: no printed word, symbol, hypothesis or number of the counted statement or of its proof was altered. Layout and class: the article's own class is \texttt{amsart} with packages including amsfonts, color, amsthm, amsmath, amscd, inputenc, t1enc, eucal, indentfirst, graphicx, graphics, pict2e, epic, url; the wrapper is the lane's pinned \texttt{amsart} editorial wrapper at 10pt on A4, which restates the theorem-like environments the retained text needs and adds the article's own macro definitions that the retained text needs (\texttt{\textbackslash wT}), with extra wrapper packages (bm, bbm, dsfont, xspace, cancel, mathtools). The delivered numbering is the unit's own. Assets: no retained span contains a figure environment, a graphics-inclusion command, a PDF-page inclusion, a table or a TikZ picture; no image file is copied into this package, the package declares no asset dependency and no image file is redistributed with it. Licence: version arXiv:2502.19196v2 of this article is distributed under the Creative Commons Attribution 4.0 International licence, as recorded in the arXiv licence metadata for this exact version and shown on the abstract page https://arxiv.org/abs/2502.19196v2. A full rights notice is delivered at licenses/arxiv-2502.19196-CC-BY-4.0.txt. Characters: every retained excerpt is pure ASCII, so the delivered engine, XeLaTeX with its default Latin Modern text font, reports no missing character.
\begin{equation} \label{eq-vlsz}
\wT_H(x,y)=\mathbb{E}\left[\prod_{v\in V(H)}Z_v(\underline{x})\right]=\int_{[0,1]^{V(H)}}.\prod_{v\in V(H)}Z_v(\underline{x})\ d\underline{x}.
\end{equation}

\begin{Th} We have
$$\lim_{n\to \infty}\wT_{H_{n,n,n}}(x,0)^{1/n}=\max_{s\in [0,1]}\left(s+x(1-s)\right)\left(xs+(1-x)\frac{s^2}{2}\right)=\begin{cases}\frac{1}{3\sqrt{3}}\frac{x^3}{x-1} &\text{if}\ \ x\geqslant\sqrt{3}, \\ 
\frac{1}{2}(x+1) &\text{if}\ \ 1< x\leqslant   \sqrt{3},\end{cases}$$
and 
$$\lim_{n\to \infty}\wT_{H_{n,n,n}}(0,x)^{1/n}=\max_{t\in [0,1]} t\left(\frac{t^2}{2}+\left(\frac{1}{2}-\frac{t^2}{2}\right)x\right)=\begin{cases}\frac{1}{3\sqrt{3}}\frac{x^{3/2}}{(x-1)^{1/2}} &\text{if}\ \ x\geqslant\frac{3}{2},  \\ 
\frac{1}{2} &\text{if}\ \ 1< x\leqslant   \frac{3}{2}.
\end{cases}$$
In particular, if $x\geqslant\sqrt{3}$ we have
$$\lim_{n\to \infty}\left(\wT_{H_{n,n,n}}(x,0)\wT_{H_{n,n,n}}(0,x)\right)^{1/n}=\left(\frac{x^3}{9(x-1)}\right)^{3/2}.$$
\end{Th}

\begin{proof}[Sketch of the proof.]
As before we generate the random permutation on \\ $V(H)=(A\cup C)\cup B$ by first generating an $x_v\in (0,1)$ uniformly at random for all $v\in V$, and then we take the relative order of $x_v$'s. Let
$$t=\max_{v\in A}x_v\ \ \ \text{and}\ \ \ s=\max_{v\in B}x_v.$$
Let $v_A$ and $v_B$ be the vertices, where these maximums are achieved.
(Note that $t$ is only maximum on $A$ and not on $A\cup C$.) As before we use equation~\ref{eq-vlsz}:
\begin{equation*} 
\wT_H(x,y)=\mathbb{E}\left[\prod_{v\in V(H)}Z_v(\underline{x})\right]=\int_{[0,1]^{V(H)}}.\prod_{v\in V(H)}Z_v(\underline{x})\ d\underline{x}.
\end{equation*}
and we aim to compress this integral using $s$ and $t$.

If $s<t$, then the contribution of this case to $\wT_{H_{n,n,n}}(x,0)$ is
\begin{equation} \label{int1}
I_1:=n^2\int_0^1\int_s^1(s+x(t-s))^{n-1}\left(\int_{0}^s(s_i+x(1-s_i))\, ds_i\right)^{n-1}(s+x(1-s))\, dt ds.
\end{equation}
Here $s+x(t-s)$ is the contribution of a vertex $w\in A$ that is not $v_A$. If $x_w=s_i<s$ for some $w\in B$ that is not $v_B$, then the attached leaf in $C$ has contribution $s_i+x(1-s_i)$. (Note that attached leaf vertex $u$ can have a value bigger than $t$ as $t$ was only a maximum among the vertices of $A$.) The last term $s+x(1-s)$ is simply the contribution of the leaf attached to $v_B$. 

If $t<s$, then the contribution of this case to $\wT_{H_{n,n,n}}(x,0)$ is
\begin{equation} \label{int2}
I_2:=n^2\int_0^1\int_s^1t^{n-1}\left(\int_{0}^t(s_i+x(1-s_i))\, ds_i+\int_{t}^sx(1-s_i)\, ds_i\right)^{n-1}x(1-s)\, dtds.
\end{equation} 
Note that we can assume that none of the vertices of $B$ are active since the contribution of such a term would be $0$.
Here, the contribution of a vertex $w\in A\setminus \{v_A\}$ is simply $t$ as it can never be active since $t<s$. If we have a vertex $w\in B$ that is not $v_B$ with $x_w=s_i<s$, then we need to distinguish two cases. If $s_i<t$, then $w$ is not active, and so the attached leaf has a contribution $s_i+x(1-s_i)$. If $s_i>t$ then, in order to make $w$ inactive the attached leaf $u$ should have a value $x_u>s_i$, but then it is automatically active and its contribution is $x(1-s_i)$. The contribution of the leaf attached to $v_B$ is $x(1-s)$ as it has to be active to make $v_B$ inactive.

We have
\begin{align*}
\lim_{n\to \infty}I_1^{1/n}&=\max_{0\leqslant s\leqslant   t\leqslant 1}(s+x(t-s))\left(xs+(1-x)\frac{s^2}{2}\right)\\
&=\max_{0\leqslant s\leqslant 1}(s+x(1-s))\left(xs+(1-x)\frac{s^2}{2}\right)
\end{align*}
The function $(s+x(1-s))\left(xs+(1-x)\frac{s^2}{2}\right)$ has either a maximum at $\left(1-\frac{1}{\sqrt{3}}\right)\frac{x}{x-1}$ if this is less than $1$, or at $s=1$. In the first case, its value is $\frac{1}{3\sqrt{3}}\frac{x^3}{x-1}$. In the second case,  its value is $\frac{1}{2}(x+1)$.

On the other hand,
$$\lim_{n\to \infty}I_2^{1/n}=\max_{0\leqslant   t\leqslant   s\leqslant   1}t\left((1-x)\frac{t^2}{2}+xt+x(s-t)-x\left(\frac{s^2}{2}-\frac{t^2}{2}\right)\right)=
\max_{0\leqslant   t\leqslant   s\leqslant   1} t\left(xs-x\frac{s^2}{2}+\frac{t^2}{2}\right)$$
Observe that at $s=t=1$ its value is $\frac{1}{2}(x+1)$. On the other hand, this is the maximum as  dropping the condition $t\leqslant s$ would lead to the following chain of inequalities
\begin{align*}
\lim_{n\to \infty}I_2^{1/n}&=\max_{0\leqslant t\leqslant s\leqslant 1} t\left(xs-x\frac{s^2}{2}+\frac{t^2}{2}\right)\\
&\leqslant \max_{0\leqslant t, s\leqslant 1} t\left(xs-x\frac{s^2}{2}+\frac{t^2}{2}\right)\\
&=\max_{0\leqslant s\leqslant 1}\left(x\left(s-\frac{s^2}{2}\right)+\frac{1}{2}\right)\\
&=\frac{1}{2}(x+1).
\end{align*}
So $\lim_{n\to \infty}I_2^{1/n}=\frac{1}{2}(x+1)$ which is always at most $\lim_{n\to \infty}I_1^{1/n}$. Then we have
\begin{align*}
\lim_{n\to \infty}\wT_{H_{n,n,n}}(x,0)^{1/n}&=\lim_{n\to \infty}(I_1+I_2)^{1/n}\\
&=\max\left(\lim_{n\to \infty}I_1^{1/n},\lim_{n\to \infty}I_2^{1/n}\right)\\
&=\lim_{n\to \infty}I_1^{1/n}.
\end{align*}
This proves the first statement. 

The statement concerning $\wT_{H_{n,n,n}}(0,x)$ is actually simpler. 
As before, let
$$t=\max_{v\in A}x_v\ \ \ \text{and}\ \ \ s=\max_{v\in B}x_v,$$
and let $v_A$ and $v_B$ be the vertices, where these maximums are achieved. Observe that we only have to consider the case $t\leqslant s$, otherwise the vertex $v_A$ is active and $X_{v_A}=0$. If $t\leqslant s$, then for any $v\in A$ different from $v_A$ the value $x_v$ can be anything between $[0,t]$ without $v$ being active. If for a vertex $v\in B$ different from $v_B$ we have $x_v=s_i$, then its unique neighbor in $C$ cannot take bigger value than $s_i$, otherwise it becomes active and the contribution of this permutation would be $0$. If $s_i<t$, then $v$ is not active, while if $s_i\in (s,t)$, then $v$ is active.  Putting all these together, we get that
\begin{align*}
\wT_{H_{n,n,n}}(0,x)&=n^2\int_0^1\int_0^st^{n-1}\left(\int_0^ts_i\, ds_i+x\int_t^ss_i\, ds_i\right)^{n-1}s\, dtds\\
&=n^2\int_0^1\int_0^st^{n-1}\left(\frac{t^2}{2}+x\left(\frac{s^2}{2}-\frac{t^2}{2}\right)\right)^{n-1}s\, dt ds.
\end{align*}
From this we get that
\begin{align*}
\lim_{n\to \infty}\wT_{H_{n,n,n}}(0,x)^{1/n}&=\max_{0\leqslant   t\leqslant   s\leqslant   1}t\left(\frac{t^2}{2}+x\left(\frac{s^2}{2}-\frac{t^2}{2}\right)\right)\\
&=\max_{0\leqslant   t\leqslant   1}t\left(\frac{t^2}{2}+x\left(\frac{1}{2}-\frac{t^2}{2}\right)\right)
\end{align*}
This expression has a maximum either at $t=\left(\frac{x}{3(x-1)}\right)^{1/2}$ if $x\geqslant\frac{3}{2}$ or at $1$ if $x\leqslant   \frac{3}{2}$. In the first case, its value is $\frac{1}{3\sqrt{3}}\frac{x^{3/2}}{(x-1)^{1/2}}$. In the second case, its value is $\frac{1}{2}$. Since $\wT_H(1,1)=1$ for any bipartite graph $H$ by definition, this completes the sketch of the proof.
\end{proof}

\end{document}
